\documentclass[11pt]{article}
\usepackage[margin=1in]{geometry}
\usepackage{amsmath,amssymb,amsthm}
\usepackage{array}
\usepackage{parskip}
\usepackage{booktabs}

\newtheorem{theorem}{Theorem}
\newtheorem{lemma}[theorem]{Lemma}
\newtheorem{corollary}[theorem]{Corollary}
\newtheorem{proposition}[theorem]{Proposition}
\theoremstyle{definition}
\newtheorem{definition}[theorem]{Definition}
\theoremstyle{remark}
\newtheorem{remark}[theorem]{Remark}

\newcommand{\Z}{\mathbb{Z}}
\newcommand{\R}{\mathbb{R}}
\newcommand{\chiG}{\chi}

\title{Disproof of conjecture \texttt{00000001556}}
\author{earthking11}
\date{2026-09-14}

\begin{document}
\maketitle

\section*{Verdict: FALSE}

We disprove conjecture \texttt{00000001556} as stated. The refutation is
unconditional and elementary: the explicit map
\[
  c_5(x,y) \;=\; (x + 2y) \bmod 5
\]
is a proper $5$-colouring of the lattice distance graph for \emph{both} standard
readings of the distance set $D=\{1,2,4\}$ simultaneously. Hence the chromatic
number is at most $5$, and the claimed value $7$ is false. In fact the exact
values are $5$ (squared-norm reading) and $3$ (Euclidean-norm reading).

\section{The conjecture, quoted verbatim}

From \texttt{conjectures/00000001556.md}:

\begin{quote}
\textbf{English.} Conjecture: The chromatic number of the planar distance graph
$G(\Z^2, D)$ with $D = \{1, 2, 4\}$ is $7$ (the exact chromatic number of a
fixed small distance set).
\end{quote}

The conjecture is filed in both English and Chinese. The Chinese version is
reproduced verbatim (in Unicode) in \texttt{README.md}; the two versions are
identical in content. The PDF font does not embed CJK glyphs, so the Chinese is
kept in the markdown copy rather than duplicated here.

\section{Definitions and the definitional ambiguity}

\begin{definition}[Distance graph]
$G(\Z^2, D)$ has vertex set $\Z^2$, and two lattice points $u,v$ are adjacent
iff the difference $u-v$ realises a ``distance'' in $D$.
\end{definition}

The conjecture is written with three scalars $D=\{1,2,4\}$, so $D$ is a set of
distances, not of integer vectors. There are two standard and equally natural
readings:

\begin{itemize}
\item \textbf{Squared-norm reading.} Lattice-colouring work often states a
  distance set as squared norms, because squared distances are integral on the
  lattice. Here $(dx,dy)$ is a displacement iff
  \[
    dx^2 + dy^2 \in \{1,2,4\}.
  \]
  This gives the $12$ vectors $(\pm1,0),(0,\pm1),(\pm1,\pm1),(0,\pm2),(\pm2,0)$.
\item \textbf{Euclidean-norm reading.} The literal reading: $(dx,dy)$ is a
  displacement iff
  \[
    \sqrt{dx^2+dy^2} \in \{1,2,4\},
  \]
  i.e.\ $dx^2+dy^2 \in \{1,4,16\}$. This gives the $12$ axis vectors
  \[
    (\pm1,0),\ (0,\pm1),\ (\pm2,0),\ (0,\pm2),\ (\pm4,0),\ (0,\pm4).
  \]
\end{itemize}

The union of the two displacement sets has $16$ vectors:
\[
\begin{array}{ll}
D_{\mathrm{sq}} = \{(\pm1,0),(0,\pm1),(\pm1,\pm1),(\pm2,0),(0,\pm2)\}, & |D_{\mathrm{sq}}|=12,\\[2pt]
D_{\mathrm{euc}} = \{(\pm1,0),(0,\pm1),(\pm2,0),(0,\pm2),(\pm4,0),(0,\pm4)\}, & |D_{\mathrm{euc}}|=12,\\[2pt]
D_{\mathrm{union}} = D_{\mathrm{sq}} \cup D_{\mathrm{euc}}, & |D_{\mathrm{union}}|=16.
\end{array}
\]
The four extra vectors are $(\pm4,0),(0,\pm4)$.

\begin{remark}[``Planar'']
Here ``planar distance graph'' names a graph drawn on the integer lattice, not a
planar graph in the graph-theoretic sense; $G(\Z^2,D)$ is far from planar.
\end{remark}

\section{The witness colouring}

\begin{theorem}\label{thm:main}
Define $c_5 \colon \Z^2 \to \{0,1,2,3,4\}$ by $c_5(x,y) = (x+2y) \bmod 5$. Then
$c_5$ is a proper colouring of the union graph $G(\Z^2, D_{\mathrm{union}})$, and
therefore of $G(\Z^2, D_{\mathrm{sq}})$ and $G(\Z^2, D_{\mathrm{euc}})$ as well.
Consequently
\[
  \chiG\bigl(G(\Z^2,D)\bigr) \le 5 < 7
\]
under both standard readings, and the conjectured value $7$ is false.
\end{theorem}

\begin{proof}
Suppose $u=(x,y)$ and $u+d = (x+dx, y+dy)$ are adjacent for some
$d=(dx,dy) \in D_{\mathrm{union}}$. Then
\[
  c_5(u+d) - c_5(u) \;\equiv\; (x+dx) + 2(y+dy) - (x+2y)
  \;\equiv\; dx + 2\,dy \pmod 5 .
\]
The displacements of the union and the corresponding residues $dx+2dy \bmod 5$
are listed in the table below; none is $0 \pmod 5$. Hence $c_5(u+d) \neq c_5(u)$
for every adjacent pair, so $c_5$ is proper. Since it uses $5$ colours,
$\chiG \le 5$. Every edge of either reading is an edge of the union graph, so
the same colouring is proper for each reading.
\end{proof}

\begin{center}
\begin{tabular}{c c r r r l}
\toprule
displacement $d=(dx,dy)$ & reading(s) & $dx+2dy$ & $\bmod\ 5$
  & $dx^2+dy^2$ & Euclidean length \\
\midrule
$(1,0)$   & both    & $1$  & $1$ & $1$  & $1$ \\
$(-1,0)$  & both    & $-1$ & $4$ & $1$  & $1$ \\
$(0,1)$   & both    & $2$  & $2$ & $1$  & $1$ \\
$(0,-1)$  & both    & $-2$ & $3$ & $1$  & $1$ \\
$(1,1)$   & squared & $3$  & $3$ & $2$  & $\sqrt2$ \\
$(1,-1)$  & squared & $-1$ & $4$ & $2$  & $\sqrt2$ \\
$(-1,1)$  & squared & $1$  & $1$ & $2$  & $\sqrt2$ \\
$(-1,-1)$ & squared & $-3$ & $2$ & $2$  & $\sqrt2$ \\
$(2,0)$   & both    & $2$  & $2$ & $4$  & $2$ \\
$(-2,0)$  & both    & $-2$ & $3$ & $4$  & $2$ \\
$(0,2)$   & both    & $4$  & $4$ & $4$  & $2$ \\
$(0,-2)$  & both    & $-4$ & $1$ & $4$  & $2$ \\
$(4,0)$   & euclid  & $4$  & $4$ & $16$ & $4$ \\
$(-4,0)$  & euclid  & $-4$ & $1$ & $16$ & $4$ \\
$(0,4)$   & euclid  & $8$  & $3$ & $16$ & $4$ \\
$(0,-4)$  & euclid  & $-8$ & $2$ & $16$ & $4$ \\
\bottomrule
\end{tabular}
\end{center}

All sixteen residues are nonzero modulo $5$, which proves Theorem~\ref{thm:main}.

\begin{remark}[Periodicity and the de Bruijn--Erd\H{o}s argument]
The colouring $c_5$ is defined by a formula on all of $\Z^2$, so it needs no
compactness argument to be global. Nevertheless, the general principle
underlying lattice colourings is worth recording: by the de Bruijn--Erd\H{o}s
compactness theorem, the chromatic number of a (locally finite) graph is the
supremum of the chromatic numbers of its finite subgraphs. For
$G(\Z^2,D)$ this says $\chiG = \sup_{F \subset \Z^2 \text{ finite}} \chiG(F)$,
so a finite clique already certifies a lower bound and a periodic colouring
certifies a global upper bound. The two bounds below exploit exactly this.
\end{remark}

\section{Exact values, with cliques as lower bounds}

\subsection*{Squared-norm reading: $\chiG = 5$}

\emph{Lower bound.} The five points
\[
  C_5 = \{(0,0),\,(1,0),\,(-1,0),\,(0,1),\,(0,-1)\}
\]
form a clique: the ten pairwise differences are
$(\pm1,0),(0,\pm1),(0,\pm2),(\pm2,0),(\pm1,\pm1),(\pm1,\mp1)$, all of squared
norm $1$, $2$ or $4$. Hence $\chiG \ge 5$. (This is the centre of the lattice
together with its four nearest neighbours.)

\emph{Upper bound.} Theorem~\ref{thm:main} gives $\chiG \le 5$ via $c_5$.
Therefore $\chiG = 5$ exactly.

\subsection*{Euclidean-norm reading: $\chiG = 3$}

\emph{Lower bound.} The three points
\[
  C_3 = \{(0,0),\,(1,0),\,(2,0)\}
\]
form a clique: the pairwise differences are $(1,0)$, $(2,0)$, $(1,0)$ of
Euclidean length $1$, $2$, $1$, all in $\{1,2,4\}$. Hence $\chiG \ge 3$.

\emph{Upper bound.} The map
\[
  c_3(x,y) = (x+y) \bmod 3
\]
is proper: for a displacement $d=(dx,dy) \in D_{\mathrm{euc}}$ we have
$dx+dy \equiv \pm1,\pm2,\pm4 \not\equiv 0 \pmod 3$. Therefore $\chiG \le 3$, and
$\chiG = 3$ exactly.

\subsection*{Union graph: $\chiG = 5$}

The clique $C_5$ lies in the squared reading, hence in the union graph, giving
$\chiG \ge 5$; the colouring $c_5$ gives $\chiG \le 5$. So the union graph also
has chromatic number exactly $5$.

\begin{corollary}
Under either standard reading of $D=\{1,2,4\}$, the conjectured value $7$ is
incorrect; the true values are $5$ (squared) and $3$ (Euclidean).
\end{corollary}

\section{Caveats: Manhattan and the continuous plane}

\subsection*{A non-standard Manhattan reading gives $\chiG \ge 9 > 7$}

If instead $D$ is read as a set of \emph{Manhattan} (taxicab) distances
$|dx|+|dy| \in \{1,2,4\}$, the displacement set is strictly larger than the
union above (it adds $(\pm1,\pm3),(\pm3,\pm1),(\pm2,\pm2)$). The nine points
\[
  \{(0,0),\,(1,\pm1),\,(2,0),\,(2,\pm2),\,(3,\pm1),\,(4,0)\}
\]
are pairwise at Manhattan distance $1$, $2$ or $4$, so they form a $9$-clique
and $\chiG \ge 9$. The colouring $c_5$ is \emph{not} proper here, since the
displacement $(3,1)$ has $3+2\cdot1=5\equiv0 \pmod 5$. Thus $7$ is false under
this reading too.

\subsection*{The continuous plane $\R^2$ is a different problem}

The only reading that would rescue the value $7$ is to reinterpret the vertex
set $\Z^2$ as the continuous plane $\R^2$. There, the chromatic number of the
plane with forbidden distance set $\{1,2,4\}$ is \emph{not} determined by the
elementary argument above; the Hadwiger--Nelson problem leaves the chromatic
number of the plane in $\{5,6,7\}$, so $7$ cannot be excluded in the continuum.
This is not the conjecture as filed: the conjecture explicitly writes
$G(\Z^2,D)$, the integer lattice, and for the lattice the refutation stands.
A referee should read the present note as disproving the lattice statement, not
the (open) continuum question.

\section{Reproducibility}

All claims above are checked by \texttt{reproduce.py} (Python~3 standard library
only):

\begin{quote}\ttfamily
\$ python3 reproduce.py\\
... (displacement sets for both readings, the 16-row residue table,\\
brute-force properness of $c_5$ on $[-200,200]^2$ and of $c_3$ on\\
$[-100,100]^2$, the $5$-clique and $3$-clique, the Manhattan $9$-clique)\\
PASS: all checks verified; conjecture 00000001556 is FALSE.
\end{quote}

It prints \texttt{PASS} and exits $0$ exactly when every check holds, and exits
non-zero otherwise.

The LaTeX document is compiled with
\begin{quote}\ttfamily
tectonic --outdir build main.tex
\end{quote}
producing \texttt{build/main.pdf}.

\section{Formalisation}

The refutation is formalised in the Lean~4 project under \texttt{lean4/}
(core Lean only, \texttt{import Std}, no Mathlib, no \texttt{sorry}). The
displacement list \texttt{D} is the $16$-vector union of both readings, and the
colourings are
\begin{quote}\ttfamily
def color  (p : Int \(\times\) Int) : Int := (p.1 + 2 * p.2) \% 5\\
def color3 (p : Int \(\times\) Int) : Int := (p.1 + p.2) \% 3
\end{quote}
The main statements are
\begin{quote}\ttfamily
theorem D\_ok :\\
\hspace*{2em}D.all (fun e => decide ((e.1 + 2*e.2) \% 5 \(\ne\) 0)) = true\\
theorem color\_shift\_ne (x y : Int) (d : Int \(\times\) Int)\\
\hspace*{2em}(hd : d \(\in\) D) :\\
\hspace*{4em}color (x + d.1, y + d.2) \(\ne\) color (x, y)\\
theorem C5\_clique :\\
\hspace*{2em}C5.all (fun p => C5.all (fun q => pairAdj p q)) = true\\
theorem Deuc\_ok :\\
\hspace*{2em}Deuc.all (fun e => decide ((e.1 + e.2) \% 3 \(\ne\) 0)) = true\\
theorem color3\_shift\_ne (x y : Int) (d : Int \(\times\) Int)\\
\hspace*{2em}(hd : d \(\in\) Deuc) :\\
\hspace*{4em}color3 (x + d.1, y + d.2) \(\ne\) color3 (x, y)\\
theorem C3\_clique :\\
\hspace*{2em}C3.all (fun p => C3.all (fun q => pairAdj p q)) = true
\end{quote}
where \texttt{pairAdj p q} is \texttt{p == q || D.contains (q.1 - p.1, q.2 - p.2)}.
The residue checks and the clique certificates are pure \texttt{decide}
computations; the properness theorems extract the nonzero residue from
\texttt{D\_ok} via \texttt{List.all\_eq\_true} and close with \texttt{omega}.
See \texttt{lean4/README.md} for the statement table, the axiom audit, and the
scope note.

\section*{Status against the submission rules}

Rule 3 requires a LaTeX source, a PDF document, and a Lean~4 project.

\begin{itemize}
\item \textbf{LaTeX source} --- \texttt{main.tex}, a standalone \texttt{article}
  using \texttt{geometry}, \texttt{amsmath}, \texttt{amssymb}, \texttt{amsthm},
  \texttt{array}, \texttt{parskip}, \texttt{booktabs}; compiles with
  \texttt{tectonic main.tex}.
\item \textbf{PDF} --- at \texttt{build/main.pdf}, produced by
  \texttt{tectonic --outdir build main.tex}.
\item \textbf{Lean~4 project} --- \texttt{lean4/} with \texttt{lean-toolchain}
  pinned to \texttt{leanprover/lean4:v4.33.1}, \texttt{lakefile.toml}
  (library \texttt{Main}, no dependencies), \texttt{Main.lean},
  \texttt{Check.lean} (\texttt{\#print axioms}), and \texttt{lean4/README.md}.
  \texttt{lake build} exits $0$ and the audit reports no \texttt{sorryAx}.
\end{itemize}

\end{document}
